\documentclass[11pt]{article}

\setlength{\parindent}{0pt}
\usepackage{xltxtra}
\usepackage{hyperref}
\hypersetup{hidelinks}
\usepackage{url}
\urlstyle{tt}
\usepackage{xcolor}
\definecolor{CVBlue}{RGB}{23,110,191}
\usepackage{calc}
\usepackage{graphicx}
\usepackage{tikz}
\usetikzlibrary{calc}
\usepackage{fontspec}
\usepackage{xeCJK}
\usepackage{enumitem}
\CJKsetecglue{}

%% 主文档字体设置
\setmainfont[
    Path = fonts/Main/,
    Extension = .otf,
    BoldFont = texgyretermes-bold.otf,
]{texgyretermes-regular.otf}

% 中文字体设置
\setCJKmainfont[
    Path = fonts/hansans/,
    Extension = .ttf,
    BoldFont = NotoSansSC-Bold.ttf,
]{NotoSansSC-Regular.otf}

\usepackage{relsize}
\usepackage{xspace}
\usepackage{fontawesome}
\usepackage[
    a4paper,
    left=1.2cm,
    right=1.2cm,
    top=1.5cm,
    bottom=1cm,
    nohead
]{geometry}

\renewcommand{\baselinestretch}{1.3} 
\usepackage{titlesec}
\setlist{noitemsep}
\setlist[itemize]{topsep=0.25em, leftmargin=*}
\setlist[enumerate]{topsep=0.25em, leftmargin=*}

\newlength{\interProjectSpacing}
\setlength{\interProjectSpacing}{0.7em}
\newcommand{\projectsep}{\vspace{\interProjectSpacing}}
\newlength{\intraProjectTitleSep}
\setlength{\intraProjectTitleSep}{0.3em}
\newcommand{\titlebreak}{\\[\intraProjectTitleSep]}
\newlength{\intraProjectListTopSep}
\setlength{\intraProjectListTopSep}{0.1em}

\titleformat{\section}
{\large\bfseries\raggedright}
{}{0em}
{}
[{\color{CVBlue}\titlerule}]
\titlespacing*{\section}{0cm}{*1.2}{*1.0}

\begin{document}
\pagenumbering{gobble}

\centerline{\LARGE\bfseries{徐焱}} 

\centerline{\normalsize{\faPhone\ +49 01626043616 \quad \faEnvelopeO\ \href{mailto:xu1398970260@163.com}{xu1398970260@163.com}}} 

\section{\makebox[\widthof{\faGraduationCap}][c]{\color{CVBlue}\faGraduationCap}\ 教育背景}    
\textbf{德国弗莱堡大学 (University of Freiburg)} \hfill 2025.10 -- 至今\\[0.3em]
嵌入式系统工程 (ESE) \quad 硕士研究生
\begin{itemize}
    \item 核心课程：微系统技术与工艺(MST)、模拟电路设计基础、半导体器件物理、集成电路设计
\end{itemize}
\textbf{武汉轻工大学} \hfill 2020.09 -- 2024.06\\[0.3em]
通信工程 \quad 本科
\begin{itemize}
    \item 核心课程：电路分析、模拟电子技术、数字电子技术、信号与系统
\end{itemize}

\section{\makebox[\widthof{\faCogs}][c]{\color{CVBlue}\faCogs}\ 专业技能}
\begin{itemize}
    \item \textbf{模拟/混合信号：} 掌握运放、滤波器、ADC/DAC设计原理；熟悉信号完整性与阻抗匹配。
    \item \textbf{EDA工具：} 熟练使用 \textbf{Multisim}；了解 \textbf{Cadence Virtuoso}、\textbf{Hspice/Spectre} 仿真流程。
    \item \textbf{软硬底层：} 精通C/C++，熟悉ARM架构与底层驱动；熟练使用示波器、万用表及信号发生器。
    \item \textbf{系统环境：} 熟悉Linux环境，具备Shell自动化脚本编写与Git版本控制能力。
\end{itemize}

\section{\makebox[\widthof{\faUsers}][c]{\color{CVBlue}\faUsers}\ 项目经历}

\textbf{高精度生理信号模拟前端(AFE)采集系统研发} \hfill \titlebreak
项目描述：主导生理信号模拟前端链路设计，实现微伏级信号采集与处理。
\begin{itemize}[topsep=\intraProjectListTopSep]
    \item \textbf{链路设计：} 设计基于运算放大器的多级调理电路，实现信号的高精度放大与预处理。
    \item \textbf{仿真验证：} 利用 \textbf{Multisim} 完成50Hz陷波与低通滤波链路仿真，验证频率响应特性。
    \item \textbf{系统集成：} 通过抗混叠滤波与阻抗匹配优化，完成了从传感器到ADC的完整信号链闭环验证。
\end{itemize}

\projectsep

\textbf{车载以太网通信可靠性研究} \hfill \titlebreak
项目描述：针对TSN车载场景，开展物理层信号分析与性能优化。
\begin{itemize}[topsep=\intraProjectListTopSep]
    \item \textbf{信号分析：} 使用示波器监测车载以太网物理层电气信号，保障高频传输下的电气可靠性。
    \item \textbf{性能优化：} 通过自动化脚本评估传输抖动与延迟，优化网络吞吐量 20\%。
\end{itemize}

\section{\makebox[\widthof{\faList}][c]{\color{CVBlue}\faList}\ 获奖情况 \& 其他}
\begin{itemize}
    \item 中国高校计算机大赛-网络技术挑战赛 省部级二等奖 / 三等奖
    \item \textbf{语言能力：} Native Mandarin, C1 English, A1 German (GRE 329, TOEFL 95)
\end{itemize}

\end{document}
